\documentclass[aspectratio=169]{beamer}
\usepackage{amsmath}
\usepackage{amssymb}
\usepackage{xcolor}

% ================================================================
% Quick Questions deck.
%
% One question per slide, no answers: the teacher jumps to a topic
% from the contents slide and picks questions for mini whiteboards.
% Within each topic the ten questions get gradually harder.
% ================================================================

\setbeamertemplate{navigation symbols}{}
\setbeamertemplate{footline}{}

% Topic divider.
\newcommand{\qtopic}[1]{%
  \section{#1}%
  \begin{frame}[plain]
    \vfill
    \begin{center}
      {\usebeamercolor[fg]{structure}\Huge\bfseries #1}
    \end{center}
    \vfill
  \end{frame}%
}

% A question slide: the question on its own, centred, nothing else on
% the slide. The optional argument drops the type size for the wordier
% questions, so nothing runs off the slide.
\newcommand{\qq}[2][\Huge]{%
  \begin{frame}
    \vfill
    \begin{center}
      \begin{minipage}{0.92\textwidth}
        \centering #1 #2
      \end{minipage}
    \end{center}
    \vfill
  \end{frame}%
}

% Contents-slide bullets, colour-coded by difficulty band:
% green for the first four topics, orange for the next four,
% red for the last three.
\setbeamertemplate{section in toc}{%
  \ifnum\inserttocsectionnumber<5
    \textcolor{green!55!black}{$\bullet$}%
  \else\ifnum\inserttocsectionnumber<9
    \textcolor{orange}{$\bullet$}%
  \else
    \textcolor{red}{$\bullet$}%
  \fi\fi
  ~\inserttocsection%
}

\title{Fractions: Quick Questions}
\subtitle{Mini whiteboard questions}
\author{}
\date{}

\begin{document}

\frame{\titlepage}

\begin{frame}{Contents}
  \small
  \tableofcontents
\end{frame}

%================================================================
\qtopic{Simplifying fractions}
%================================================================

\qq{Simplify $\dfrac{4}{8}$}
\qq{Simplify $\dfrac{6}{9}$}
\qq{Simplify $\dfrac{12}{20}$}
\qq{Simplify $\dfrac{24}{36}$}
\qq{Simplify $\dfrac{45}{60}$}
\qq{Simplify $\dfrac{42}{70}$}
\qq{Simplify $\dfrac{84}{144}$}
\qq{Simplify $\dfrac{91}{117}$}
\qq{Simplify $\dfrac{315}{420}$}
\qq[\LARGE]{Simplify $\dfrac{1001}{1309}$\\[0.4em]
\normalsize (Hint: both numbers are multiples of $7$ and of $11$.)}

%================================================================
\qtopic{Ordering fractions}
%================================================================

\qq[\LARGE]{Put these in order, smallest first:\\[0.6em]
$\dfrac{1}{4}\qquad \dfrac{1}{2}\qquad \dfrac{1}{8}$}

\qq[\LARGE]{Put these in order, smallest first:\\[0.6em]
$\dfrac{3}{7}\qquad \dfrac{6}{7}\qquad \dfrac{2}{7}$}

\qq[\LARGE]{Put these in order, smallest first:\\[0.6em]
$\dfrac{1}{2}\qquad \dfrac{3}{4}\qquad \dfrac{5}{8}$}

\qq[\LARGE]{Put these in order, smallest first:\\[0.6em]
$\dfrac{2}{3}\qquad \dfrac{5}{6}\qquad \dfrac{3}{4}$}

\qq[\LARGE]{Put these in order, smallest first:\\[0.6em]
$\dfrac{3}{5}\qquad \dfrac{7}{10}\qquad \dfrac{13}{20}$}

\qq[\LARGE]{Put these in order, smallest first:\\[0.6em]
$\dfrac{5}{9}\qquad \dfrac{2}{3}\qquad \dfrac{7}{12}$}

\qq[\LARGE]{Put these in order, \textbf{largest} first:\\[0.6em]
$\dfrac{4}{5}\qquad \dfrac{7}{9}\qquad \dfrac{11}{15}$}

\qq[\LARGE]{Put these in order, smallest first:\\[0.6em]
$\dfrac{5}{8}\qquad \dfrac{7}{11}\qquad \dfrac{9}{14}$}

\qq[\LARGE]{Which is closer to $1$: $\dfrac{8}{9}$ or $\dfrac{11}{12}$?\\[0.4em]
\normalsize How do you know?}

\qq[\LARGE]{Write a fraction that lies between\\[0.6em]
$\dfrac{5}{8}$ \quad and \quad $\dfrac{2}{3}$}

%================================================================
\qtopic{Mixed numbers and improper fractions}
%================================================================

\qq{Write $1\dfrac{1}{2}$ as an improper fraction}
\qq{Write $2\dfrac{3}{4}$ as an improper fraction}
\qq{Write $\dfrac{7}{2}$ as a mixed number}
\qq{Write $4\dfrac{2}{5}$ as an improper fraction}
\qq{Write $\dfrac{23}{4}$ as a mixed number}
\qq{Write $5\dfrac{7}{8}$ as an improper fraction}
\qq{Write $\dfrac{100}{7}$ as a mixed number}
\qq{Write $12\dfrac{5}{9}$ as an improper fraction}
\qq{Write $-3\dfrac{2}{5}$ as an improper fraction}
\qq[\LARGE]{The mixed number $6\dfrac{a}{7}$ is the same as $\dfrac{45}{7}$.\\[0.6em]
Find $a$.}

%================================================================
\qtopic{Adding and subtracting fractions}
%================================================================

\qq{$\dfrac{1}{2}+\dfrac{1}{4}$}
\qq{$\dfrac{2}{3}+\dfrac{1}{6}$}
\qq{$\dfrac{3}{4}-\dfrac{1}{8}$}
\qq{$\dfrac{1}{3}+\dfrac{1}{4}$}
\qq{$\dfrac{3}{5}-\dfrac{1}{4}$}
\qq{$\dfrac{5}{6}+\dfrac{3}{8}$}
\qq{$\dfrac{7}{10}-\dfrac{2}{15}$}
\qq{$\dfrac{2}{3}+\dfrac{3}{4}+\dfrac{1}{6}$}
\qq{$\dfrac{5}{12}-\dfrac{7}{18}$}
\qq[\LARGE]{What must be added to $\dfrac{5}{9}$\\[0.4em]
to make $\dfrac{7}{8}$?}

%================================================================
\qtopic{Converting fractions, decimals and percentages}
%================================================================

\qq[\LARGE]{Write $\dfrac{1}{2}$ as a decimal and as a percentage}
\qq[\LARGE]{Write $0.25$ as a fraction and as a percentage}
\qq[\LARGE]{Write $30\%$ as a decimal and as a fraction in its simplest form}
\qq[\LARGE]{Write $\dfrac{3}{4}$ as a percentage, and $0.6$ as a percentage}
\qq[\LARGE]{Write $\dfrac{7}{20}$ as a decimal and as a percentage}
\qq[\LARGE]{Write $0.08$ as a percentage and as a fraction in its simplest form}
\qq[\LARGE]{Write $\dfrac{3}{8}$ as a decimal and as a percentage}
\qq[\LARGE]{Write $12.5\%$ as a fraction in its simplest form}
\qq[\LARGE]{Write $\dfrac{5}{6}$ as a decimal, and as a percentage to $1$ decimal place}
\qq[\LARGE]{Write $137.5\%$ as a mixed number in its simplest form}

%================================================================
\qtopic{Ordering fractions, decimals and percentages}
%================================================================

\qq[\LARGE]{Put these in order, smallest first:\\[0.6em]
$0.5\qquad \dfrac{1}{4}\qquad 60\%$}

\qq[\LARGE]{Put these in order, smallest first:\\[0.6em]
$0.7\qquad \dfrac{3}{4}\qquad 70\%$}

\qq[\LARGE]{Put these in order, smallest first:\\[0.6em]
$\dfrac{2}{5}\qquad 45\%\qquad 0.35$}

\qq[\LARGE]{Put these in order, smallest first:\\[0.6em]
$0.62\qquad \dfrac{5}{8}\qquad 63\%$}

\qq[\LARGE]{Put these in order, \textbf{largest} first:\\[0.6em]
$\dfrac{7}{8}\qquad 0.87\qquad 88\%$}

\qq[\LARGE]{Put these in order, smallest first:\\[0.6em]
$\dfrac{1}{3}\qquad 0.3\qquad 33\%$}

\qq[\LARGE]{Put these in order, smallest first:\\[0.6em]
$0.045\qquad 4.5\%\qquad \dfrac{1}{20}$}

\qq[\LARGE]{Put these in order, smallest first:\\[0.6em]
$\dfrac{5}{6}\qquad 83\%\qquad 0.835$}

\qq[\LARGE]{Put these in order, smallest first:\\[0.6em]
$\dfrac{2}{7}\qquad 0.28\qquad 29\%$}

\qq[\LARGE]{Put these in order, smallest first:\\[0.6em]
$\dfrac{3}{8}\qquad 0.38\qquad 37.5\%\qquad \dfrac{3}{7}$\\[0.5em]
\normalsize What do you notice?}

%================================================================
\qtopic{Fractions of amounts (no calculator)}
%================================================================

\qq{$\dfrac{1}{2}$ of $40$}
\qq{$\dfrac{1}{4}$ of $60$}
\qq{$\dfrac{1}{5}$ of $45$}
\qq{$\dfrac{3}{4}$ of $40$}
\qq{$\dfrac{2}{3}$ of $36$}
\qq{$\dfrac{3}{5}$ of $45$}
\qq{$\dfrac{5}{8}$ of $72$}
\qq{$\dfrac{2}{7}$ of $350$}
\qq[\LARGE]{$\dfrac{5}{6}$ of a number is $45$.\\[0.4em]
What is the number?}
\qq[\LARGE]{$\dfrac{3}{8}$ of $\dfrac{2}{5}$ of $200$}

%================================================================
\qtopic{Fractions of amounts (calculator)}
%================================================================

\qq{$\dfrac{1}{4}$ of $386$}
\qq{$\dfrac{3}{5}$ of $245$}
\qq{$\dfrac{2}{7}$ of $483$}
\qq{$\dfrac{3}{16}$ of $1024$}
\qq{$\dfrac{5}{9}$ of $738$}
\qq[\LARGE]{$\dfrac{7}{8}$ of $\pounds 312.40$}
\qq{$\dfrac{11}{13}$ of $2405$}
\qq[\LARGE]{$\dfrac{4}{7}$ of $3.5$ kg.\\[0.4em]
Give your answer in grams.}
\qq[\LARGE]{$\dfrac{5}{8}$ of a number is $215$.\\[0.4em]
What is the number?}
\qq[\LARGE]{$\dfrac{7}{12}$ of $4$ hours $48$ minutes.\\[0.4em]
Give your answer in hours and minutes.}

%================================================================
\qtopic{Four operations with mixed numbers and negatives}
%================================================================

\qq{$1\dfrac{1}{2}+2\dfrac{1}{4}$}
\qq{$3\dfrac{4}{5}-1\dfrac{2}{5}$}
\qq{$2\dfrac{1}{2}+\left(-1\dfrac{1}{4}\right)$}
\qq{$4\dfrac{1}{6}-2\dfrac{3}{4}$}
\qq{$2\dfrac{1}{2}\times 4$}
\qq{$-3\dfrac{1}{3}+1\dfrac{5}{6}$}
\qq{$1\dfrac{1}{2}\times 2\dfrac{2}{3}$}
\qq{$3\dfrac{3}{4}\div 1\dfrac{1}{2}$}
\qq{$-5\dfrac{1}{3}\div \dfrac{2}{3}$}
\qq[\LARGE]{$-2\dfrac{1}{4}\times\left(-1\dfrac{1}{3}\right)+1\dfrac{1}{2}$}

%================================================================
\qtopic{Word problems with mixed numbers}
%================================================================

\qq[\LARGE]{A recipe needs $1\dfrac{1}{2}$ cups of flour.\\[0.4em]
Alex makes a double batch.\\[0.4em]
How much flour is needed?}

\qq[\LARGE]{A plank is $3\dfrac{1}{4}$\,m long.\\[0.4em]
A piece $1\dfrac{1}{2}$\,m long is cut off.\\[0.4em]
How much plank is left?}

\qq[\LARGE]{Priya has two ribbons, $2\dfrac{1}{3}$\,m and $1\dfrac{5}{6}$\,m long.\\[0.4em]
What is their total length?}

\qq[\LARGE]{Sam runs $2\dfrac{3}{4}$ miles on Monday and $3\dfrac{1}{2}$ miles on Tuesday.\\[0.4em]
How much further did he run on Tuesday?}

\qq[\LARGE]{A tin holds $4\dfrac{1}{2}$ litres of paint.\\[0.4em]
Five jugs of $\dfrac{3}{4}$ of a litre are poured out.\\[0.4em]
How much paint is left?}

\qq[\LARGE]{A wall is $7\dfrac{1}{2}$\,m wide.\\[0.4em]
Panels are $1\dfrac{1}{4}$\,m wide.\\[0.4em]
How many panels fit across the wall?}

\qq[\LARGE]{A bottle holds $1\dfrac{1}{3}$ litres of water.\\[0.4em]
Jo drinks $\dfrac{3}{8}$ of it.\\[0.4em]
How much water does Jo drink?}

\qq[\LARGE]{A journey takes $2\dfrac{1}{4}$ hours.\\[0.4em]
Maria has driven $\dfrac{2}{3}$ of the way.\\[0.4em]
How long is left, in minutes?}

\qq[\LARGE]{A rope $12\dfrac{1}{2}$\,m long is cut into pieces $1\dfrac{3}{4}$\,m long.\\[0.4em]
How many whole pieces are there, and how much rope is left over?}

\qq[\large]{A recipe for $4$ people uses $2\dfrac{1}{2}$ cups of rice and $1\dfrac{3}{4}$ litres of stock.\\[0.5em]
How much rice and stock are needed for $10$ people?}

%================================================================
\qtopic{Recurring decimals to fractions}
%================================================================

\qq[\LARGE]{Use an algebraic method to write\\[0.5em]
$0.\dot{3}$ \quad as a fraction}

\qq[\LARGE]{Use an algebraic method to write\\[0.5em]
$0.\dot{7}$ \quad as a fraction}

\qq[\LARGE]{Use an algebraic method to write\\[0.5em]
$0.\dot{4}\dot{5}$ \quad as a fraction}

\qq[\LARGE]{Use an algebraic method to write\\[0.5em]
$0.\dot{1}\dot{2}$ \quad as a fraction in its simplest form}

\qq[\LARGE]{Use an algebraic method to write\\[0.5em]
$0.\dot{2}\dot{7}$ \quad as a fraction in its simplest form}

\qq[\LARGE]{Use an algebraic method to write\\[0.5em]
$0.\dot{1}2\dot{3}$ \quad as a fraction in its simplest form}

\qq[\LARGE]{Use an algebraic method to write\\[0.5em]
$0.1\dot{6}$ \quad as a fraction in its simplest form}

\qq[\LARGE]{Use an algebraic method to write\\[0.5em]
$0.8\dot{3}$ \quad as a fraction in its simplest form}

\qq[\LARGE]{Use an algebraic method to write\\[0.5em]
$1.\dot{2}\dot{7}$ \quad as a mixed number in its simplest form}

\qq[\LARGE]{Use an algebraic method to show that\\[0.5em]
$0.\dot{9}=1$}

\end{document}
